\chapter[Hiperparâmetros]{Hiperparâmetros}
\label{ap:hiperparametros}

Este apêndice consolida os hiperparâmetros das etapas já executadas do pipeline.
O Quadro~\ref{tab:hp_gerais} reúne os parâmetros comuns à geração de pares
candidatos e à avaliação; o Quadro~\ref{quad:hp_b1} detalha o \emph{baseline} de
regras (B1) e o Quadro~\ref{quad:hp_b2} o \emph{baseline} de regressão logística
(B2).

\begin{quadro}[htbp]
  \centering
  \caption{Parâmetros gerais do pipeline.}
  \label{tab:hp_gerais}
  \begin{tabular}{ll}
  \toprule
  Parâmetro & Valor \\
  \midrule
  Semente aleatória (padrão) & 42 \\
  Sementes p/ análise de variância & 42, 123, 7, 2024, 31337 \\
  Janela máxima de pares (\texttt{max\_gap}) & 200 caracteres \\
  Partição treino/dev/teste & 80/10/10 (nível de documento) \\
  \emph{Bootstrap} (reamostragens) & 1000 (método percentil, IC95\%) \\
  \bottomrule
  \end{tabular}\\[4pt]
  {\footnotesize Fonte: elaborado pelo autor.}
\end{quadro}

\begin{quadro}[htbp]
  \centering
  \caption{Hiperparâmetros do \emph{baseline} de regras (B1).}
  \label{quad:hp_b1}
  \begin{tabular}{ll}
  \toprule
  Parâmetro & Valor \\
  \midrule
  Contexto à esquerda (R1) & 30 caracteres \\
  Limiar de distância (R2) & 50 caracteres \\
  Rótulo padrão (R3) & \texttt{no\_relation} \\
  \bottomrule
  \end{tabular}\\[4pt]
  {\footnotesize Fonte: elaborado pelo autor (\texttt{scripts/baselines/baseline\_rules.py}).}
\end{quadro}

\begin{quadro}[htbp]
  \centering
  \caption{Hiperparâmetros do \emph{baseline} de regressão logística (B2).}
  \label{quad:hp_b2}
  \begin{tabular}{ll}
  \toprule
  Parâmetro & Valor \\
  \midrule
  \emph{TF-IDF} (entre/esquerda) & n-gramas 1--2, \emph{lowercase}, \texttt{min\_df}=2 \\
  Vocabulário máximo (entre) & 50000 \\
  Faixas de distância (\emph{buckets}) & 0, 10, 25, 50, 100, 200 \\
  Reescala de classes & \texttt{class\_weight="balanced"} \\
  Regularização (C) & 1.0 \\
  \emph{Solver} & \texttt{lbfgs} (padrão do \emph{scikit-learn}) \\
  Máximo de iterações & 1000 \\
  Características totais & 94257 \\
  Candidatos (treino) & 831902 \\
  \bottomrule
  \end{tabular}\\[4pt]
  {\footnotesize Fonte: elaborado pelo autor (\texttt{scripts/baselines/baseline\_lr.py} e \texttt{experiments/results/baseline\_lr.json}).}
\end{quadro}

O Quadro~\ref{quad:hp_biobertpt} reúne os hiperparâmetros \emph{planejados} para o
modelo principal (BioBERTpt), conforme \texttt{docs/plano\_notebook\_biobertpt.md}.
Os valores definitivos serão confirmados após a execução do treino.

\begin{quadro}[htbp]
  \centering
  \caption{Hiperparâmetros planejados do modelo proposto (BioBERTpt).}
  \label{quad:hp_biobertpt}
  \begin{tabular}{ll}
  \toprule
  Parâmetro & Valor planejado \\
  \midrule
  \emph{Encoder} & \texttt{pucpr/biobertpt-all} \\
  Representação de entrada & marcadores de entidade tipados (\texttt{[E1]}/\texttt{[E2]}) \\
  \emph{Pooling} & concatenação dos estados de início ($h_{E1}, h_{E2}$) \\
  Cabeçote & 1 camada linear (1536 $\rightarrow$ 3), \emph{dropout} 0,1 \\
  Perda & \emph{cross-entropy} ponderada (\texttt{class\_weight="balanced"}) \\
  \emph{max\_length} & 256 \emph{tokens} \\
  Taxa de aprendizado & 2e-5 (AdamW) \\
  \emph{Batch} (treino) & 16 + acúmulo 2 (efetivo 32) \\
  \emph{Batch} (avaliação) & 64 \\
  Épocas (máx.) & 5 (com \emph{early stopping}, paciência 2) \\
  \emph{warmup\_ratio} & 0,1 \\
  \emph{weight\_decay} & 0,01 \\
  \emph{fp16} & sim \\
  Seleção de modelo & macro-F1 no \emph{dev} \\
  Sementes de treino & 42, 123, 1337 \\
  \emph{Bootstrap} & 1000, semente 42 (fixa) \\
  Hardware & Google Colab, \emph{Tesla} T4 \\
  \bottomrule
  \end{tabular}\\[4pt]
  {\footnotesize Fonte: elaborado pelo autor (\texttt{docs/plano\_notebook\_biobertpt.md}).}
\end{quadro}

Os valores efetivamente utilizados substituirão os planejados após a execução do
treino do modelo.
